% ============================================================================
%  Part IV -- From selection to synthesis: guided diffusion  (section fragment)
%  \input-ready; uses ONLY macros from preamble.tex. No preamble here.
%  Cross-part refs used: \eqref{eq:gradxK} (Part I kernel input-gradient),
%  \eqref{eq:Ustd} \eqref{eq:Uda} (Part III utilities),
%  \ref{subsec:numerics} (Part III utility numerics / truncation ceiling).
%  Notation clash discipline (preamble): the diffusion schedule \bsched/\aret/\abar
%  and timestep \dt are DISTINCT from the prior half-width \bwid of the localized
%  region and the locality mask \loc; \dt is a scalar diffusion-time index, NOT a
%  grid coordinate -- the prior is over the 2-D grid only.
% ============================================================================

\section{Guided diffusion for input synthesis}\label{sec:diffusion}

Parts I--III chose the next input by \emph{selecting} the highest-utility member
of a fixed pool. This part develops the natural extension: \emph{synthesize} a new
input, drawn from the input (data) distribution $\px$, that maximizes the same
Gaussian-process (GP) utility for a given target model. A denoising diffusion model
supplies a learned prior over $\px$, and the GP utility gradient---the very object
of Part~III---steers generation toward high-utility inputs. The two ingredients
divide the labor cleanly: the utility gradient (smooth, low-pass;
\S\ref{sec:diff-motiv}) fixes \emph{where} in input space to move, and the
diffusion score supplies the high-frequency structure of $\px$ that a utility
gradient can never produce on its own.

% ----------------------------------------------------------------------------
\subsection{Motivation: synthesize rather than select}\label{sec:diff-motiv}

\paragraph{Why direct gradient ascent on the utility fails.}
Suppose we optimize an input $x^{*}$ directly to maximize a GP utility $U(x^{*})$
by gradient ascent. The gradient flows through the arc-cosine kernel's
input-gradient $\nabla_{x}\Kf$ (Part~I, \eqref{eq:gradxK}), which carries the
structured (grid) prior $\Cmat$ and hence its Gaussian smoothing block
$\Csm$---a convolution with a Gaussian kernel, i.e.\ a \emph{low-pass filter}. Two
consequences follow, independent of the inducing points, the training data, or
which utility is used:
\begin{enumerate}[leftmargin=1.4em,itemsep=1pt,topsep=2pt]
  \item the utility gradient is \emph{inherently smooth}, so ascent iterates lack
        the high-frequency structure (sharp transitions, fine texture, local
        contrast) that characterizes draws from $\px$; and
  \item ascent pushes grid points outside the bounded input box. The arc-cosine
        self-kernel is positively homogeneous, $\Kf(x,x)=x\transpose
        \Cmat x+\sigz^{2}=\lVert x\rVert_{\Cmat}^{2}+\sigz^{2}$, so larger norms
        are rewarded and the utility diverges under input scaling (the
        norm-scaling divergence of Part~III).
\end{enumerate}
Such iterates are not admissible inputs: to lie in the support of $\px$ they must
be indistinguishable from genuine draws from the input distribution.

\paragraph{Why a subspace constraint is not enough.}
Restricting the optimization to a low-dimensional subspace (PCA space, or the
$\Cmat$-eigenspace of Part~III) does not cure this. PCA space fixes only the
\emph{second-order} statistics; under approximate stationarity of $\px$ its
eigenvectors tend to Fourier modes (Szeg\H{o}'s theorem), so PCA-space ascent is
essentially Fourier low-pass filtering. Neither PCA space nor the
$\Cmat$-eigenspace enforces the higher-order structure (phase coherence, sparse
gradients) that places a sample in the support of $\px$: a subspace constraint is a
convergence aid, not a $\px$-fidelity constraint.

\paragraph{The decomposition, and why diffusion.}
The synthesis task has two objectives that pull apart cleanly:
(i)~\emph{informativeness}---$x^{*}$ maximizes the utility (reduces output
uncertainty); and (ii)~\emph{fidelity}---$x^{*}$ is a plausible draw from the input
distribution $\px$. Gradient ascent handles (i) but not (ii). A diffusion model
trained on samples from $\px$ learns the \emph{score}
$\nabla_{x}\log p_{\dt}(x)$ at every noise level $\dt$; at low noise this encodes
the full statistical structure of $\px$. Guided generation adds the two:
\begin{equation}
\underbrace{\score(\xt,\dt)}_{\text{high-freq.\ structure}}
\;+\;
\underbrace{\gw\,\nabla_{\xt}U(\xhat)}_{\text{low-freq.\ utility bias}} .
\label{eq:diff-decomp}
\end{equation}

\begin{remark}[Division of labor]
The smoothness that ruined direct optimization becomes acceptable under
\eqref{eq:diff-decomp}: the utility only has to \emph{bias} generation toward
informative regions, not prescribe individual grid points. The score model supplies
all orders of the statistics of $\px$ (trained once, reused across targets and
utilities); the utility gradient supplies the target-specific, low-frequency
direction. They act at complementary spatial-frequency scales.
\end{remark}

\paragraph{Connection to the selection setting.}
In the fixed-pool active-learning loop, one \emph{selects} the pool input of
highest expected information gain about the latent function. Diffusion-guided
synthesis is the extension: \emph{manufacture} an input of even higher utility that
still lies in the support of $\px$---closing the gap between ``best available pool
input'' and ``best attainable informative input.'' This part is framed by three
questions: (1)~can we generate inputs that are simultaneously plausible under $\px$
and high-utility?  (2)~how does that input change as the GP sharpens from
$\Ntr=50$ to $300$ training inputs?  (3)~how much utility does the $\px$-fidelity
constraint cost---the gap between the unconstrained pure-utility optimum and the
diffusion-guided one, as a function of guidance strength $\gw$?

% ----------------------------------------------------------------------------
\subsection{Diffusion fundamentals}\label{sec:diff-fund}

\subsubsection{Forward (noising) process}
Let $\xz\in\Reals^{d}$ be a clean input ($d=64\times64=4096$ here). Fix a variance
schedule $\bsched\in(0,1)$, $\dt=1,\dots,T$, and define the per-step signal
retention $\aret=1-\bsched$ and its cumulative product
$\abar=\prod_{s=1}^{\dt}\alpha_{s}=\prod_{s=1}^{\dt}(1-\beta_{s})$. The forward
process adds Gaussian noise one step at a time,
\begin{equation}
q(\xt\mid x_{\dt-1})=\Gauss\!\big(\xt;\ \sqrt{1-\bsched}\,x_{\dt-1},\ \bsched\,I\big),
\label{eq:forward-step}
\end{equation}
and admits a closed-form marginal at arbitrary $\dt$ (the key property---no
intermediate steps are needed):
\begin{equation}
q(\xt\mid\xz)=\Gauss\!\big(\xt;\ \sqrt{\abar}\,\xz,\ (1-\abar)\,I\big).
\label{eq:forward-diff}
\end{equation}
Equivalently, by the reparameterization
\begin{equation}
\xt=\sqrt{\abar}\,\xz+\sqrt{1-\abar}\,\epsilon,\qquad \epsilon\sim\Gauss(0,I).
\label{eq:reparam}
\end{equation}
The schedule is chosen so that $\bar\alpha_{T}\approx0$, hence
$q(x_{T}\mid\xz)\approx\Gauss(0,I)$: at the end of the forward process all
information about $\xz$ is destroyed.

\subsubsection{Score $\leftrightarrow$ noise, and the training objective}
Differentiating the log of the forward marginal \eqref{eq:forward-diff} gives the
score in closed form, which identifies it with the added noise:
\begin{equation}
\nabla_{\xt}\log q(\xt\mid\xz)
   =-\frac{\xt-\sqrt{\abar}\,\xz}{1-\abar}
   =-\frac{\epsilon}{\sqrt{1-\abar}}
\quad\Longrightarrow\quad
\nabla_{\xt}\log p_{\dt}(\xt)\approx\score(\xt,\dt)=-\frac{\epsth(\xt,\dt)}{\sqrt{1-\abar}}.
\label{eq:score-noise}
\end{equation}
A network trained to predict the noise therefore approximates the score. Training
minimizes the simplified denoising-score-matching (DSM) $L_{2}$ loss
\begin{equation}
\mathcal{L}_{\mathrm{DSM}}
  =\E_{\dt\sim U(1,T)}\,\E_{\xz\sim\px}\,\E_{\epsilon\sim\Gauss(0,I)}
   \Big[\big\lVert\epsth(\xt,\dt)-\epsilon\big\rVert^{2}\Big],
\label{eq:dsm}
\end{equation}
with $\xt$ formed by \eqref{eq:reparam}. Each mini-batch draws a fresh
$(\xz,\dt,\epsilon)$; the network learns all noise levels though it sees each
infrequently.

\subsubsection{Reverse (denoising) process and the Tweedie estimate}
The generative reverse process is Gaussian,
$p_{\theta}(x_{\dt-1}\mid\xt)=\Gauss(x_{\dt-1};\mu_{\theta}(\xt,\dt),\sigma_{\dt}^{2}I)$,
with mean read off from the noise predictor
\begin{equation}
\mu_{\theta}(\xt,\dt)=\frac{1}{\sqrt{1-\bsched}}
   \Big[\,\xt-\frac{\bsched}{\sqrt{1-\abar}}\,\epsth(\xt,\dt)\Big],
\label{eq:reverse-mean}
\end{equation}
and fixed variance $\sigma_{\dt}^{2}$, either the simple choice
$\sigma_{\dt}^{2}=\bsched$ or the posterior form
$\tilde\beta_{\dt}=\bsched(1-\bar\alpha_{\dt-1})/(1-\abar)$; in practice
$\sigma_{\dt}^{2}=\bsched$ works well. Sampling runs $\dt=T,\dots,1$ from
$x_{T}\sim\Gauss(0,I)$, drawing $x_{\dt-1}=\mu_{\theta}+\sigma_{\dt}z$ with
$z\sim\Gauss(0,I)$ ($z=0$ at $\dt=1$). At any $\dt$ the noise prediction yields the
\emph{Tweedie estimate} of the clean input---the model's posterior mean
$\E[\xz\mid\xt]$:
\begin{equation}
\xhat=\frac{\xt-\sqrt{1-\abar}\,\epsth(\xt,\dt)}{\sqrt{\abar}}.
\label{eq:tweedie}
\end{equation}
It is coarse at high noise and increasingly detailed as $\dt\downarrow$. The
Tweedie estimate \eqref{eq:tweedie} is the workhorse of guidance
(\S\ref{sec:diff-guide}): it maps a noisy intermediate to a clean input on which
the GP utility can be evaluated, and it is a differentiable function of $\xt$.

\begin{remark}[$T-1$ start]
At $\dt=T$ the schedules used here have $\bar\alpha_{T}\!\sim\!10^{-3}$, so
$1/\sqrt{\bar\alpha_{T}}\!\sim\!31$ and the first reverse step amplifies any
prediction error $\sim\!31\times$, causing divergence. Sampling therefore starts
the loop at $\dt=T-1$ (where $1/\sqrt{\alpha}\!\approx\!2$, stable).
\end{remark}

\subsubsection{The U-Net noise predictor and DDIM acceleration}
Because $\epsth(\xt,\dt)$ maps a noisy input to a same-size tensor, it is
parameterized by an encoder--decoder with \emph{skip connections} (a U-Net): a
plain bottleneck would discard exactly the fine detail that distinguishes a draw
from $\px$ from noise, whereas skip connections carry per-resolution detail
directly to the decoder. The scalar timestep $\dt$ enters through a sinusoidal
embedding (Transformer-style),
\[
\mathrm{emb}(\dt)_{2k}=\sin\!\big(\dt/10000^{2k/d_{\mathrm{emb}}}\big),
\qquad
\mathrm{emb}(\dt)_{2k+1}=\cos\!\big(\dt/10000^{2k/d_{\mathrm{emb}}}\big),
\]
passed through a small MLP and \emph{added} (broadcast over space) to each block's
feature map---a global ``how much to activate at this noise level'' modulation.

Full DDPM sampling needs $T-1=999$ sequential U-Net passes. \emph{DDIM} replaces
the stochastic reverse chain by a deterministic one over a subsequence of $S$
timesteps $\tau_{S}>\dots>\tau_{1}$ (e.g.\ $S=50$). Going from $\dt=\tau_{i}$ to
$t'=\tau_{i-1}$,
\begin{equation}
x_{t'}=\sqrt{\bar\alpha_{t'}}\,\xhat
   +\sqrt{1-\bar\alpha_{t'}-\sigma^{2}}\;
     \frac{\xt-\sqrt{\abar}\,\xhat}{\sqrt{1-\abar}}
   +\sigma\,z,
\label{eq:ddim}
\end{equation}
with $\xhat$ from \eqref{eq:tweedie}. Here $\sigma=0$ gives a fully deterministic,
reproducible sampler; the DDPM choice of $\sigma$ recovers stochastic sampling.
Quality degrades below $S\!\approx\!20$; $S=50$ is the default.

% ----------------------------------------------------------------------------
\subsection{Guidance: steering the reverse process with the GP utility gradient}
\label{sec:diff-guide}

\subsubsection{Classifier-style guidance}
Model ``high utility'' as a label $y$ with $p(y\mid\xz)\propto\exp\!\big(\gw\,
U(\xz)\big)$, $\gw>0$. Bayes' rule splits the \emph{conditional} score into a
$\px$-fidelity term and a guidance term:
\begin{equation}
\nabla_{\xt}\log p(\xt\mid y)
  =\underbrace{\score(\xt,\dt)}_{\text{unconditional score, \eqref{eq:score-noise}}}
  +\;\gw\,\nabla_{\xt}\Ustd(\xhat(\xt,\dt)).
\label{eq:guidance}
\end{equation}
The utility is defined on \emph{clean} inputs, so it is evaluated on the Tweedie
estimate \eqref{eq:tweedie} (the standard Dhariwal--Nichol trick). Any of the
Part~III utilities may play the role of $U$ in \eqref{eq:guidance}: the standard
utility $\Ustd$ (\eqref{eq:Ustd}), shown here, or the distribution-aware utility
$\Uda$ (\eqref{eq:Uda}). Define the \emph{guidance gradient}
\begin{equation}
\ggrad=\nabla_{\xt}\Ustd\big(\xhat(\xt,\dt)\big).
\label{eq:ggrad}
\end{equation}

\subsubsection{The guided reverse update (two equivalent forms)}
Converting the guided score \eqref{eq:guidance} back to a noise prediction, and
substituting into the Tweedie map \eqref{eq:tweedie}, yields two algebraically
identical updates:
\begin{align}
\epsilon_{\mathrm{guided}} &= \epsth(\xt,\dt)-\gw\,\sqrt{1-\abar}\,\ggrad,
   \label{eq:guided-noise}\\[2pt]
\xhat^{\mathrm{guided}} &= \xhat+\frac{\gw\,(1-\abar)}{\sqrt{\abar}}\,\ggrad.
   \label{eq:guided-tweedie}
\end{align}
Form \eqref{eq:guided-noise} feeds the DDPM mean \eqref{eq:reverse-mean} in place
of $\epsth$; form \eqref{eq:guided-tweedie} feeds the DDIM update \eqref{eq:ddim}
in place of $\xhat$. The \emph{ramp factor} $(1-\abar)/\sqrt{\abar}$ in
\eqref{eq:guided-tweedie} is large at high noise ($\dt\!\approx\!T$, $\sim\!158$),
giving bold utility steering while global structure is still being decided, and
small at low noise ($\dt\!\approx\!1$, $\sim\!0.3$), preserving fine detail.

\subsubsection{The GP link: what the guidance signal is}\label{sec:diff-gplink}
The guidance gradient \eqref{eq:ggrad} is the gradient of the GP information
utility of Part~III, back-propagated through Tweedie. Both utilities are
differentiable in the input, and the gradient flows through the arc-cosine
kernel's input-gradient $\nabla_{x}\Kf$ (Part~I, \eqref{eq:gradxK}). Because
$\nabla_{x}\Kf$ carries the factor $\Cmat$---and thus the Gaussian spatial
smoother $\Csm$---the utility gradient is exactly the smooth, low-pass signal of
\S\ref{sec:diff-motiv}. This is precisely why guidance works: the utility supplies
a \emph{low-frequency} informative direction, and the score model supplies the
\emph{high-frequency} structure of $\px$ the smooth utility gradient can never
produce.

\paragraph{The locality mask $\Rightarrow$ guidance touches only the
localized-region grid points.}
The structured prior applies a Gaussian locality envelope to each grid point $j$,
\begin{equation}
\loc_{j}=\exp\!\Big(-\frac{\lVert\pixc_{j}-\rfc\rVert^{2}}{4\,\bwid^{2}}\Big),
\label{eq:rfmask}
\end{equation}
where $\pixc_{j}$ is grid point $j$'s coordinate, $\rfc$ the center of the
localized region, and $\bwid$ the \emph{kernel's localized-region half-width}.
Grid points with $\loc_{j}<10^{-3}$ are masked out (zero contribution), so
$\nabla_{x}U$ is nonzero \emph{only at localized-region grid points}. The mask area
scales as $\bwid^{2}$: at the default $\bwid=0.1$ the localized region covers
$\approx21\%$ of the $64\times64$ grid ($\approx869$ grid points), versus $5.4\%$
at $\bwid=0.05$ and $1.7\%$ at $\bwid=0.03$. The utility therefore guides only the
localized region, and \emph{the diffusion model freely generates every remaining
grid point} for $\px$-fidelity---an advantage over the L$_{2}$ approach
(\S\ref{sec:diff-approaches}), whose non-localized-region grid points are frozen at
a start input. Because $\bwid$ is a learnable kernel parameter, it must be frozen
during GP training for synthesis, or the intended mask size drifts.

\caveat{The symbol $\bwid$ in \eqref{eq:rfmask} is the \emph{kernel}
localized-region half-width (Part~I), \textbf{not} the diffusion noise-schedule
variance $\bsched$ of \eqref{eq:forward-step}; likewise the locality envelope
$\loc_{j}$ is the kernel locality mask, not the diffusion signal-retention $\aret$.
These collisions ($\beta$, $\alpha$) are the most dangerous in the whole document;
the macros keep them apart.}

\subsubsection{Full versus approximate guidance gradient}\label{sec:diff-approx}
By the chain rule $\ggrad$ factors through the utility gradient in clean-input
space and the Tweedie Jacobian,
\begin{equation}
\frac{\partial\xhat}{\partial\xt}
   =\frac{1}{\sqrt{\abar}}\Big(I-\sqrt{1-\abar}\,\frac{\partial\epsth}{\partial\xt}\Big),
\label{eq:tweedie-jac}
\end{equation}
where $\partial\epsth/\partial\xt$ is the $d\times d$ U-Net Jacobian, which is
never formed explicitly.
\begin{itemize}[leftmargin=1.4em,itemsep=1pt,topsep=2pt]
  \item \textbf{Full gradient:} retains $\partial\epsth/\partial\xt$ in
        \eqref{eq:tweedie-jac}; one U-Net forward \emph{and} one backward pass,
        $\approx2\times$ the cost of a plain step.
  \item \textbf{Approximate gradient} (default): treat the U-Net Jacobian as
        negligible. Then $\xhat$ depends on $\xt$ only through the direct
        $\xt/\sqrt{\abar}$ term, so
        \begin{equation}
        \ggrad^{\mathrm{approx}}=\frac{1}{\sqrt{\abar}}\,\nabla_{\xhat}\Ustd(\xhat).
        \label{eq:approx-grad}
        \end{equation}
        No U-Net backward pass ($\approx\tfrac12$ the cost). It is justified when
        $\epsth$ is a ``local correction'' whose Jacobian is roughly zero-mean
        (does not systematically align with or oppose the utility gradient), which
        holds best at \emph{low} noise (small $\sqrt{1-\abar}$); a hybrid
        (approximate for early $\dt>200$, full for late $\dt<200$) captures most
        of the benefit cheaply.
\end{itemize}

\subsubsection{The four approaches (A--D) and the adopted method}
\label{sec:diff-approaches}
Four guidance paradigms were catalogued; all attempt to add a fidelity term
counteracting the utility gradient's drift off the data manifold $\mathcal{M}$ (the
support of $\px$). Let $x_{c}$ denote a current input being optimized.

\begin{description}[leftmargin=1.6em,itemsep=2pt,topsep=2pt]
  \item[A --- L$_2$ penalty toward the Tweedie estimate.] Corrupt $x_{c}$ at a
        \emph{fixed} noise level $\dt^{*}$, predict $\hat\epsilon$, form $\xhat$
        (\eqref{eq:tweedie}), hold it fixed (no gradient through it), and penalize
        the distance to it:
        \begin{equation}
        \mathcal{L}_{\mathrm{A}}=-\Uda(x_{c})+\frac{\greg}{2}\,
           \big\lVert x_{c}-\xhat\big\rVert^{2},
        \qquad
        x_{c}\leftarrow x_{c}+\eta\,\nabla\Uda+\eta\,\greg\,(\xhat-x_{c}).
        \label{eq:approachA}
        \end{equation}
        Here $\greg$ is the approach-A regularization weight. Since
        $x_{c}-\xhat\propto(\hat\epsilon-\epsilon)$, \eqref{eq:approachA} is a
        Score-Distillation-Sampling gradient with a different time weighting.
        \emph{Weaknesses:} no signal near $\mathcal{M}$ (for $x_{c}$ already on
        $\mathcal{M}$, $\xhat\!\approx\!x_{c}$ so $\hat\epsilon-\epsilon$ is just
        reconstruction noise---the penalty is reactive, not preventive); the
        Tweedie estimate of an amplified input is itself amplified, so norm growth
        is not prevented; $\dt^{*}$ is arbitrary; a single fixed noise sample is
        biased.
  \item[B --- direct score.] Use $\score(\xt,\dt^{*})=-\hat\epsilon/\sqrt{1-
        \bar\alpha_{\dt^{*}}}$ directly, chain-ruled to $x_{c}$-space. Uses
        $\hat\epsilon$ alone (not $\hat\epsilon-\epsilon$), so it has signal on
        $\mathcal{M}$ but high variance, and estimates the score of the
        \emph{blurred} $p_{\dt^{*}}$, not $p_{0}$.
  \item[C --- score without added noise.] Feed $x_{c}$ at $\dt=1$ directly.
        Appealing ($p_{1}\!\approx\!p_{0}$) but fails: a utility-distorted $x_{c}$
        is far out-of-distribution for a network expecting ``a draw from $\px$ plus
        $\sigma\!\approx\!0.01$ noise,'' and dividing by
        $\sqrt{1-\bar\alpha_{1}}\!\approx\!0.01$ amplifies the unreliable
        prediction $\sim\!100\times$.
  \item[D --- full guided reverse diffusion (\emph{adopted}).] Do not optimize an
        existing input---\emph{generate} from noise, nudging each denoising step
        by the utility gradient via \eqref{eq:guidance}--\eqref{eq:guided-tweedie}.
        The U-Net is \emph{always} in-distribution ($\xt$ is genuinely at noise
        level $\dt$ because it came from the reverse process); generation
        traverses coarse$\to$fine over all noise levels; there is no $\dt^{*}$ to
        choose; the output is a proper sample from the guided distribution. Costs:
        the U-Net forward (with gradients in full mode), $T$ (or $\sim\!50$ DDIM)
        sequential steps, utility evaluated on rough early Tweedie estimates,
        stochastic output (best-of-$K$), and no start-point control.
\end{description}

\begin{remark}[Why the score is a genuine restoring force]
Naive per-coordinate clipping $\Pi_{\mathcal{B}}$ (each $x_{i}$ clamped to
$[x_{\min},x_{\max}]$) acts on each grid point independently and ignores the joint
distribution $\px$: it lands on the boundary of the input box but far from
$\mathcal{M}$. In the guided score \eqref{eq:guidance} the score term is instead a
\emph{structured} restoring force on the joint distribution---if a region saturates
and loses its characteristic spectrum, $p_{\dt}(\xt)$ drops sharply and the score
rearranges grid points to restore the correlations of $\px$ while accommodating the
utility bias. The limitation is out-of-distribution failure: $\epsth$ was trained
only on the forward trajectory (samples from $\px$ plus Gaussian noise), so too
large a $\gw$ drives $\xt$ where $\epsth$ extrapolates and the restoring force
becomes unreliable. Thus $\gw$ must be large enough to explore high-utility regions
yet small enough to keep $\xt$ where the score approximation is valid.
\end{remark}

\subsubsection{The rate guard}
At high noise the Tweedie estimate is coarse and can produce extreme grid points,
hence extreme rates that fall in the regime where the truncated utility is
inaccurate (Part~III, \S\ref{subsec:numerics}). The generator therefore zeroes the
guidance gradient whenever the predicted log-rate is too high:
\begin{equation}
\ggrad\leftarrow 0
\quad\text{whenever}\quad
\exp(\mug)\ge f_{\max},
\qquad
\mug=\gain\,\pmean+\latz,
\label{eq:fmax}
\end{equation}
with $f_{\max}=100$ and $\mug$ the predicted log-rate mean (gain $\gain$ times the
posterior latent mean $\pmean$ plus the offset $\latz$). This guard activates mostly in
the first $\sim\!10$ steps (high $\dt$) and keeps generation in the low-rate regime
where the (truncated) utility matches a Monte-Carlo evaluation. Because the utility
landscape over random seeds is sparse (most seeds land at $U\!\approx\!0$),
generation is run best-of-$K$: retain the highest-utility input over $K$ random
seeds.

% ----------------------------------------------------------------------------
\subsection{Open issues}\label{sec:diff-status}

Two remaining questions are mathematical rather than procedural.

\paragraph{Utility-accuracy ceiling at high rate.}
Both utilities inherit the Poisson-sum truncation at $\rmax$ analyzed in Part~III
(\S\ref{subsec:numerics}): above rate $\approx\!50$ the standard utility $\Ustd$
overflows to $+\infty$ (its analytic noise-entropy term is unbounded), while the
distribution-aware utility $\Uda$ stays finite but is silently wrong (its marginal
entropy collapses toward $0$). Guidance is therefore reliable only in the low-rate
regime, rate $\approx\!0.01$--$1$, where both utilities agree with a Monte-Carlo
evaluation to within $\sim\!3\%$; the rate guard \eqref{eq:fmax} is exactly what
confines generation to that regime. Extending accurate utility evaluation above the
truncation bound---so that guidance remains valid at high rate---is open.

\paragraph{Single-sample distribution-aware estimator.}
In the guided reverse path the distribution-aware utility \eqref{eq:Uda} is
evaluated with a single conditioning sample ($n_{\mathrm{mc}}=1$, the current
input) rather than a Monte-Carlo average over $\px$; enlarging $n_{\mathrm{mc}}$ is
a straightforward but untested refinement.
